% !TeX root = ../main.tex
% ---------------------------------------------------------------------------
% Segmental Lining Design
% Scope extension: not derived from the RSE3010 weekly material.
% Key sources: Muir Wood (1975); Duddeck & Erdmann (1985); Einstein & Schwartz
% (1979); Janssen joint model; ITA WG2 Guidelines for the Design of Segmental
% Tunnel Linings (2019); ACI 544.7R; fib Bulletin 83.
% ---------------------------------------------------------------------------
\chapter{Segmental Lining Design}
\label{ch:segmental-linings}


\section{Ring function, geometry and materials}
\label{sec:ring-function-geometry-and-materials}

\subsection{Function and configuration of a segmental ring}
\label{sec:ring-function}
\scopenote{The ring as simultaneously the TBM reaction frame, the primary
support and the permanent lining --- three roles with conflicting design
drivers.}

\subsection{Ring geometry and segmentation}
\label{sec:ring-geometry}
\scopenote{Universal and left/right tapered rings, key segment arrangement,
segment count, slenderness, and how curve radius and build tolerance set the
taper.}

\subsection{Materials}
\label{sec:segment-materials}
\scopenote{Conventionally reinforced concrete, steel fibre reinforced concrete,
hybrid reinforcement, and the design code basis for each.}

\section{Load cases and ground loading}
\label{sec:load-cases-and-ground-loading}

\subsection{Load cases}
\label{sec:segment-loads}
\scopenote{Ground and water loading, grout pressure, TBM thrust jack forces,
demoulding, handling, stacking, erection, and long-term loads including
surcharge and seismic. Most segment cracking originates in the transient
cases.}

\subsection{Ground loading models}
\label{sec:ground-loading-models}
\scopenote{Full overburden, arching-reduced loading (Terzaghi), and
soil-structure interaction approaches; choosing a model consistent with
ground stiffness and construction method.}

\section{Analysis methods}
\label{sec:analysis-methods}

\subsection{Continuum solutions}
\label{sec:continuum-solutions}
\scopenote{Muir Wood, Duddeck \& Erdmann and Einstein \& Schwartz closed-form
solutions; the compressibility and flexibility ratios; full-slip and no-slip
interface assumptions.}

\subsection{Bedded beam-spring models}
\label{sec:beam-spring}
\scopenote{Radial and tangential spring stiffness, one-way ground springs, and
the sensitivity of computed moments to spring assumptions.}

\subsection{Numerical modelling of the lining}
\label{sec:lining-numerical}
\scopenote{Embedding the ring in a continuum model, staged installation, and
extracting thrust and moment for section design. Links to
Chapter~\ref{ch:numerical-modelling}.}

\section{Joint behaviour}
\label{sec:joint-behaviour}

\subsection{Radial joint behaviour}
\label{sec:radial-joints}
\scopenote{The Janssen model for a reduced contact joint, rotational
stiffness, the resulting reduction in effective ring stiffness, and Muir
Wood's effective inertia correction.}

\subsection{Circumferential joints and ring coupling}
\label{sec:circumferential-joints}
\scopenote{Ring-to-ring coupling, staggered joint arrangements, dowels and
cams, and the load sharing between adjacent rings.}

\section{Connections and watertightness}
\label{sec:connections-and-watertightness}

\subsection{Connections: bolts, dowels and guiding rods}
\label{sec:segment-connections}
\scopenote{Temporary and permanent connection functions, gasket compression
during build, and the trend toward bolt-free systems.}

\subsection{Gaskets and watertightness}
\label{sec:gaskets}
\scopenote{EPDM and hydrophilic gaskets, gap and offset performance envelopes,
required contact pressure, and specifying watertightness classes.}

\section{Local effects and build quality}
\label{sec:local-effects-and-build-quality}

\subsection{Bursting, spalling and local effects}
\label{sec:segment-local-effects}
\scopenote{Concentrated loads at joints and jack shoes, bursting
reinforcement, Leonhardt design, and the detailing that prevents corner
spalling.}

\subsection{Tolerances, build quality and ring build}
\label{sec:ring-build}
\scopenote{Ovalisation, lipping, stepping and birdsmouthing; how build quality
governs the durability and watertightness actually achieved.}

\section{Damage mechanisms and repair}
\label{sec:segment-damage}
\scopenote{Cracking patterns and their causes, acceptance criteria, and repair
methods for damaged rings.}
